\chapter[1992 --- Fischer et al.]{1992: Edmond H. Fischer, Edwin G. Krebs}
\label{ch:1992_Fischer}

\section*{Main Contribution}

\textit{Discovered reversible protein phosphorylation as a biological regulatory mechanism, explaining how cells switch proteins on and off to control metabolism.}\cite{nobel-1992,lecture-1992}

\section{Official Citation}

for their discoveries concerning reversible protein phosphorylation as a biological regulatory mechanism

\section{Background and Historical Context}

The regulation of cellular processes---metabolism, growth, differentiation, and communication---depends on the ability of cells to switch proteins on and off in response to external signals. One of a major mechanisms for achieving this molecular switching is the reversible phosphorylation of proteins: the addition of a phosphate group to specific amino acid residues and its subsequent removal. This process, elucidated through the collaborative work of Edmond H. Fischer and Edwin G. Krebs, was recognized with the 1992 Nobel Prize in Physiology or Medicine and has since been recognized as a universal regulatory mechanism governing virtually every aspect of cell biology.

Edmond Henry Fischer was born on April 6, 1920, in Shanghai, China, to French and Swiss parents. He studied chemistry at the University of Geneva and earned his doctorate in 1947 under the supervision of Robert Signer. Fischer's early research concerned the chemistry of enzymes---specifically, the structure and specificity of the active sites of proteolytic enzymes. In 1950, he moved to the University of Washington in Seattle, where he encountered Edwin Krebs and began the collaboration that would define his career.

Edwin Gerhard Krebs was born on June 6, 1918, in Lansing, Iowa. He studied medicine at Washington University in St. Louis and completed his residency in internal medicine before turning to research. After a postdoctoral fellowship with Carl Cori, who had shared the Nobel Prize in 1947 for his work on glycogen metabolism, Krebs moved to the University of Washington in 1948. His training in glycogen biochemistry provided the foundation for the discovery that would later win the Nobel Prize.

The intellectual context of their work centered on the regulation of glycogen metabolism---the process by which cells store and mobilize glucose. In the 1940s, it was known that the hormone epinephrine (adrenaline) stimulated the breakdown of glycogen in muscle and liver, while insulin promoted glycogen synthesis. The enzyme responsible for glycogen breakdown, glycogen phosphorylase, had been identified by Carl Cori and others. However, the mechanism by which hormones controlled the activity of this enzyme was unknown. A critical clue came from the observation that glycogen phosphorylase existed in two forms: an active ``a'' form and a less active ``b'' form, and that the conversion between these forms was regulated by enzymatic reactions.

\section{The Discovery/Contribution}

\subsection{The Discovery of Phosphorylase Kinase}

Fischer and Krebs's breakthrough began with a simple but powerful question: what enzyme was responsible for converting phosphorylase $b$ to phosphorylase $a$? Working with rabbit skeletal muscle---a tissue rich in glycogen and therefore rich in phosphorylase---they discovered a new enzyme, which they named phosphorylase kinase, that catalyzed the ATP-dependent phosphorylation of glycogen phosphorylase $b$, converting it to the active $a$ form.

Their initial experiments, published in a series of landmark papers in the 1950s, demonstrated that the activation of phosphorylase involved the transfer of a phosphate group from ATP to a serine residue on the phosphorylase protein. This was a startling finding because, at the time, protein phosphorylation was not recognized as a widespread regulatory mechanism. Proteins were thought to be relatively stable molecular structures, and the idea that a small chemical modification could dramatically alter a protein's activity was novel.

Fischer and Krebs went on to show that the reverse reaction---the dephosphorylation of phosphorylase $a$ back to phosphorylase $b$---was catalyzed by a separate enzyme, which they named protein phosphatase. The phosphorylation--dephosphorylation cycle thus constituted a reversible molecular switch: phosphorylation by kinase activated the enzyme, and dephosphorylation by phosphatase inactivated it.

\subsection{The Expansion to Other Proteins}

As their work progressed, Fischer and Krebs recognized that protein phosphorylation was not unique to glycogen phosphorylase. They and their colleagues systematically surveyed other metabolic enzymes and discovered that many of them were regulated by reversible phosphorylation. Glycogen synthase, the enzyme responsible for glycogen synthesis, was found to be inactivated by phosphorylation---the opposite of the effect on phosphorylase. This reciprocal regulation ensured that glycogen synthesis and breakdown were coordinated: when one pathway was active, the other was inhibited.

The discovery that protein phosphatases constituted a family of enzymes, each with its own substrate specificity, further underscored the generality of the phosphorylation--dephosphorylation mechanism. Fischer and Krebs also identified the role of calcium ions and cyclic AMP (cAMP) as second messengers that activated specific kinases, linking extracellular hormone signals to intracellular phosphorylation cascades.

The implications of these discoveries were immense. By the late 1990s, it was estimated that approximately one-third of all proteins in a eukaryotic cell are phosphorylated at any given time, and that protein kinases and phosphatases constitute roughly 2--3\% of the total protein content of a cell. The human genome encodes more than 500 protein kinases and approximately 150 protein phosphatases, reflecting the notable importance of reversible phosphorylation in cellular regulation.

\subsection{Signal Transduction Cascades}

Fischer and Krebs's work on phosphorylase kinase also anticipated the modern concept of signal transduction cascades. The activation of glycogen phosphorylase by epinephrine involves a multi-step pathway: epinephrine binds to a cell surface receptor, which activates a G protein, which stimulates adenylyl cyclase to produce cAMP, which activates protein kinase A (PKA), which phosphorylates and activates phosphorylase kinase, which phosphorylates and activates glycogen phosphorylase. Each step in this cascade amplifies the signal, allowing a single hormone molecule to trigger the mobilization of millions of glycogen molecules.

This amplification principle is a hallmark of signal transduction and is now recognized as fundamental to many physiological processes, including vision, olfaction, immune signaling, and growth factor signaling. Fischer and Krebs's elucidation of the phosphorylase activation cascade was thus one of the earliest examples of a signal transduction pathway, predating the discovery of many of its components.

\section{Impact on Modern Medicine}

The impact of Fischer and Krebs's work on modern medicine is pervasive. Protein phosphorylation is now recognized as the primary mechanism by which cells regulate their behavior in response to external signals. Dysregulation of phosphorylation pathways is implicated in a vast array of diseases, including cancer, diabetes, cardiovascular disease, neurodegenerative disorders, and autoimmune conditions.

\subsection{Cancer and Kinase Inhibitors}

Perhaps the most dramatic clinical impact has been in the field of oncology. Many oncogenes encode protein kinases that are constitutively active, driving uncontrolled cell proliferation. The development of small-molecule kinase inhibitors---including imatinib (Gleevec) for chronic myeloid leukemia, trastuzumab (Herceptin) for HER2-positive breast cancer, and erlotinib for non-small cell lung cancer---has transformed cancer treatment. These drugs, which now account for a major fraction of new anticancer agents approved by the FDA, are direct intellectual descendants of Fischer and Krebs's discovery that protein phosphorylation regulates cellular function.

\subsection{Diabetes and Metabolic Disease}

The phosphorylation of glycogen synthase and glycogen phosphorylase remains central to our understanding of glucose metabolism and insulin signaling. Defects in insulin receptor signaling---a phosphorylation cascade involving insulin receptor substrate (IRS) proteins, PI3-kinase, and Akt---are a hallmark of type 2 diabetes. Understanding these pathways has guided the development of antidiabetic drugs and continues to inform research on metabolic syndrome and obesity.

\subsection{Neurodegenerative Disease}

In the brain, protein phosphorylation plays critical roles in synaptic plasticity, memory formation, and neuronal survival. Aberrant phosphorylation of the tau protein is a hallmark of Alzheimer's disease, and the accumulation of phosphorylated tau into neurofibrillary tangles is a defining pathological feature. Understanding the kinases and phosphatases that regulate tau phosphorylation is an active area of drug development for Alzheimer's disease and related tauopathies.

\section{Legacy}

Edmond Fischer and Edwin Krebs shared the 1992 Nobel Prize in Physiology or Medicine for their discoveries. Fischer became chairman of the Department of Biochemistry at the University of Washington and continued to study the regulation of protein kinases, particularly the role of tyrosine phosphorylation in growth factor signaling. He received numerous honors, including the Albert Lasker Basic Medical Research Award in 1976. Fischer passed away on August 19, 2021, at the age of 101.

Edwin Krebs became chairman of the Department of Pharmacology at the University of Washington and later moved to the University of California, Davis, where he continued his research on signal transduction. He was also awarded the Albert Lasker Basic Medical Research Award in 1976 and received numerous other honors during his career. Krebs passed away on December 21, 2009, at the age of 91.

Their legacy endures in the countless researchers who have built upon their discoveries. The ``kinome''---the complete set of protein kinases encoded by the human genome---was mapped in the early 2000s, and large-scale proteomics projects continue to catalog the ``phosphoproteome''---the full complement of phosphorylated proteins in human cells. Fischer and Krebs's work opened a vast new field of biology and medicine that continues to yield insights and therapeutic opportunities.


\section{Modern Developments}

Fischer and Krebs' protein phosphorylation work has led to modern signal transduction research. Understanding of kinase signaling has led to targeted cancer therapies (imatinib, erlotinib, vemurafenib). Kinase inhibitors are among the most successful class of cancer drugs. Understanding of phosphatase function has revealed new therapeutic targets. CRISPR-based screens identify kinases essential for cancer cell survival, accelerating drug target discovery.


\section{Bibliography}

\begin{thebibliography}{9}
\bibitem{nobel-1992}
Nobel Prize Outreach, ``The Nobel Prize in Physiology or Medicine 1992,''\emph{NobelPrize.org}.\\
\url{https://www.nobelprize.org/prizes/medicine/1992/summary/}

\bibitem{lecture-1992}
Edmond H. Fischer, ``Nobel Lecture,''\emph{NobelPrize.org}. Winner-authored primary source; downloadable from the official Nobel Prize archive.\\
\url{https://www.nobelprize.org/prizes/medicine/1992/fischer/lecture/}
\end{thebibliography}
